% Generated draft: requires mathematical and source review.
\documentclass[11pt]{article}
\usepackage[margin=25mm]{geometry}
\usepackage{fontspec}
\setmainfont{Latin Modern Roman}
\usepackage{amsmath,amssymb,mathtools,unicode-math}
\setmathfont{Latin Modern Math}
\usepackage{enumitem,needspace,xurl,hyperref}
\setlength{\parindent}{0pt}
\setlength{\parskip}{5pt}
\setlength{\emergencystretch}{3em}
\newcommand{\N}{\mathbb{N}}
\newcommand{\Z}{\mathbb{Z}}
\newcommand{\Q}{\mathbb{Q}}
\newcommand{\R}{\mathbb{R}}
\newcommand{\C}{\mathbb{C}}
\newcommand{\F}{\mathbb{F}}
\newcommand{\A}{\mathbb{A}}
\newcommand{\PP}{\mathbb{P}}
\newcommand{\E}{\mathbb{E}}
\newcommand{\eps}{\varepsilon}
\newcommand{\dd}{\,\mathrm{d}}
\begin{document}
\section*{UnsolvedMath ID: 20001993}
\textbf{Source problem code:} catalogue entry
\subsection{Definitions and mathematical statement}
The source asks for tropical techniques to compute the chamberwise real genus-zero enumerative invariants of \(\C\mathrm P^2\) attributed to G. Tian and S. Kwon. The constraints are real Chow cycles; the invariant is described as a count of intersections of their pullbacks on the real locus of the Kontsevich stable-map moduli space. Degrees, markings, transversality, excluded discriminant and chamber conventions must be fixed before this becomes a numerical problem. The source reports that the rational nodal Severi locus embeds as a Zariski-open dense subset of the relevant stable-map space, providing the link to Mikhalkin's plane nodal-curve correspondence.

The supplied aid treats the restricted point-constraint problem for \(d\geq3\). Take \(3d-1\) distinct real points in a regular configuration \(\mathbf p\) for which the degree-\(d\) rational-curve problem is finite, transverse, nodal, irreducible and represented by birational maps, with no boundary contribution. Let \(\mathcal C\) be its connected chamber in that regular configuration locus, and let \(R_d(\mathcal C)\) denote the unsigned real curve count. This point version is not a definition of the full Chow-cycle theory in the source.

For tropical point conditions equipped with real phase data and satisfying a real correspondence theorem's hypotheses, let \(\mathcal T\) be the finite collection of contributing tropical curves and let \(m_{\R}(\Gamma)\) count their real algebraic lifts. The aid reports that
\[
A=\sum_{\Gamma\in\mathcal T}m_{\R}(\Gamma)
=R_d(\mathcal C_{\mathrm{lift}}),
\]
where \(\mathcal C_{\mathrm{lift}}\) is the chamber containing the constructed real algebraic lift of the constraints. Real signs/phases in the tropical data do not turn this into an alternating signed count: the summands here count real lifts. To certify the count for a prescribed chamber \(\mathcal C\), one must additionally certify \(\mathcal C_{\mathrm{lift}}=\mathcal C\), for example by a path in the regular locus. The aid identifies this chamber-membership step as the remaining global task.

It also reports a valuation-only obstruction: generic cubic configurations with unsigned counts 10 and 12 become identical under bare coordinatewise valuation after embedding their nonzero coordinates as constant real Puiseux series. Thus valuations alone cannot determine every chamber count. Phase data and a valid algebraic-lift/chamber certificate are material.

For \(d=3\), the aid supplies the independent numerical audit
\[
h-e=8,\qquad R_3=h+e\in\{8,10,12\},
\]
where \(h\) counts real rational nodal cubics with a nonisolated real node and \(e\) those with an isolated real node. An unsigned tropical output \(A\) is consequently required to satisfy
\[
h=\frac{A+8}{2},\qquad e=\frac{A-8}{2}.
\]
The audit is a necessary consistency check, not a substitute for chamber identification or a theorem for arbitrary real Chow constraints. Full correspondence hypotheses and detailed references for these claims are absent from the record.
\subsection{Short English statement}
Develop tropical calculations of chamberwise unsigned real curve counts in the projective plane. The supplied aid gives a phase-enriched point-constraint reduction but requires a certificate locating the algebraic lift in the desired chamber. Bare valuations and a cubic consistency audit cannot replace that step.
\subsection{Sources}
\begin{itemize}
\item Seongchun Kwon, A.11, \emph{Real Gromov--Witten invariants and tropical geometry}, AIM workshop \emph{Amoebas and tropical geometry}: \url{https://aimath.org/WWN/amoebas/amoebas.pdf}.
\item The source attributes the chamberwise invariant to G. Tian and S. Kwon and invokes Mikhalkin's plane rational nodal-curve work. No full bibliographic entries or detailed correspondence theorem are supplied; the original literature field is empty.
\item Local registry \texttt{data/unsolvedmath/problems.json}, permanent ID 20001993, code \texttt{catalogue entry}; canonical \texttt{aim-geometry-notes.json}, \texttt{notes[330]}; updated 14 August 2026.
\item Attempt 1: reduction, partially solved status, rigor label proved, novelty confidence low. The machine-generated aid is not peer reviewed and requires independent verification before citation.
\end{itemize}
\end{document}
